%-------------------------
% Currículo em LaTeX - Otimizado para ATS (Sistemas de IA de RH)
% Autor: Gustavo Luiz Heck
% Tipo: Estágio em Engenharia Mecânica / Automação
%-------------------------

\documentclass[11pt,a4paper,sans]{article}

% --- PACOTES BÁSICOS ---
\usepackage[utf8]{inputenc}
\usepackage[T1]{fontenc}
\usepackage[portuguese]{babel}
\usepackage{lmodern}

% --- AJUSTE DE MARGENS ---
\usepackage[top=1.6cm, bottom=1.6cm, left=1.8cm, right=1.8cm]{geometry}

% --- DESIGN & CORES ---
\usepackage{xcolor}
\definecolor{primaryColor}{HTML}{0F172A}    % Azul Escuro / Slate
\definecolor{accentColor}{HTML}{1D4ED8}     % Azul Royal para detalhes/seções
\definecolor{textColor}{HTML}{334155}       % Cinza Escuro para texto corrido

% --- AJUSTE DE FONTES ---
\renewcommand{\familydefault}{\sfdefault}   % Fonte Sem Serifa (Arial/Helvetica style)

% --- OUTROS PACOTES ---
\usepackage{titlesec}
\usepackage{enumitem}
\usepackage{hyperref}
\usepackage{tabularx}

% --- CONFIGURAÇÕES DE LINKS (hyperref) ---
\hypersetup{
    colorlinks=true,
    linkcolor=accentColor,
    urlcolor=accentColor,
    pdftitle={Gustavo Luiz Heck - Currículo},
    pdfauthor={Gustavo Luiz Heck}
}

% --- ESTILO DE SEÇÕES (titlesec) ---
\titleformat{\section}
  {\normalfont\large\bfseries\color{accentColor}} % Formato
  {}                                             % Label (número)
  {0pt}                                          % Espaço entre label e título
  {}                                             % Código antes
  [{\color{accentColor}\titlerule[0.8pt]}]       % Linha horizontal abaixo

\titlespacing*{\section}{0pt}{12pt}{6pt}

% --- COMANDOS CUSTOMIZADOS PARA ORGANIZAÇÃO ---
\newcommand{\experienceHeader}[4]{
    \noindent
    \begin{tabularx}{\textwidth}{@{} X r @{}}
        \textbf{\color{primaryColor}#1} | \textit{#2} & \small{#3} \\
        \small\textit{\color{textColor}#4} & \\
    \end{tabularx}
    \vspace{-2pt}
}

\newcommand{\educationItem}[4]{
    \noindent
    \begin{tabularx}{\textwidth}{@{} X r @{}}
        \textbf{\color{primaryColor}#1} & \small{#3} \\
        \small{#2} & \small\textit{#4} \\
    \end{tabularx}
    \vspace{2pt}
}

% --- RETIRAR INDENTAÇÃO PADRÃO ---
\setlength{\parindent}{0pt}

\begin{document}
\color{textColor}

% ==========================================
% CABEÇALHO (Informações Pessoais)
% ==========================================
\begin{center}
    {\Huge \textbf{\color{primaryColor}Gustavo Luiz Heck}} \\
    \vspace{6pt}
    \small Joinville/SC \ | \ (48) 99144-0068 \ | \ 
    \href{mailto:gustavoluizheck+trabalho@gmail.com}{gustavoluizheck+trabalho@gmail.com} \\
    \vspace{2pt}
    \href{https://www.linkedin.com/in/gustavo-luiz-heck-987442207/}{linkedin.com/in/gustavo-luiz-heck-987442207/}
\end{center}

\vspace{6pt}

% ==========================================
% OBJETIVO
% ==========================================
\section{Objetivo}
Estudante de Engenharia Mecânica com sólida formação técnica em Automação Industrial, em busca de \textbf{Estágio em Engenharia}. Foco em integrar conhecimentos multidisciplinares nas áreas de \textbf{Automação e Controle}, \textbf{Projetos Mecânicos} e \textbf{Engenharia de Sistemas}. Perfil adaptativo com experiência prática em liderança de equipes de desenvolvimento e gestão de projetos em empresa júnior e robótica móvel de competição.

% ==========================================
% FORMAÇÃO ACADÊMICA
% ==========================================
\section{Formação Acadêmica}

\educationItem
  {Bacharelado em Engenharia Mecânica}
  {UDESC -- Centro de Ciências Tecnológicas (CCT)}
  {2020/1 -- Presente (6ª fase em andamento)}
  {Joinville/SC}

\vspace{4pt}

\educationItem
  {Curso Técnico em Automação Industrial}
  {SENAI Tubarão}
  {Concluído em 2019/2}
  {Tubarão/SC}

% ==========================================
% EXPERIÊNCIA
% ==========================================
\section{Experiência}

\experienceHeader
  {GERM UDESC}
  {Grupo Estudantil de Robótica Móvel}
  {3 anos e 2 meses (jun/2023 -- presente)}
  {Bolsista de Extensão \& Líder (desde dez/2023) | Membro Voluntário (jun/2023 -- dez/2023)}

\begin{itemize}[nosep, leftmargin=*]
    \item \textbf{Robótica para a Inclusão Social (RISO):} Desenvolvimento e aplicação de projetos educacionais voltados à robótica móvel, circuitos eletrônicos e programação para a comunidade.
    \item \textbf{Liderança do Processo Seletivo:} Planejamento, estruturação e coordenação de estratégias de recrutamento, dinâmicas de fit cultural e avaliação técnica de novos membros.
    \item \textbf{Liderança NERO \& ROBOCOMP:} Gestão de equipes em robôs de competição. Atuação no desenvolvimento de firmware (microcontroladores Arduino) em \textbf{linguagem C/C++} e modelagem mecânica de componentes.
\end{itemize}

\vspace{6pt}

\experienceHeader
  {Konvex Jr.}
  {Empresa Júnior de Engenharia Mecânica -- UDESC-CCT}
  {1 ano e 11 meses (fev/2021 -- dez/2022)}
  {Diretor de Projetos (dez/2021 -- dez/2022) | Assessor, Consultor \& Trainee (fev/2021 -- dez/2021)}

\begin{itemize}[nosep, leftmargin=*]
    \item \textbf{Gestão de Portfólio (Diretoria):} Planejamento estratégico e controle de qualidade de projetos de consultoria mecânica, com foco no cumprimento de prazos, escopos e satisfação do cliente.
    \item \textbf{Liderança de Equipes:} Gestão direta de projetistas em soluções de engenharia, promovendo metodologias ágeis e cultura corporativa.
    \item \textbf{Modelagem 3D \& Detalhamento CAD:} Elaboração de geometrias, montagens de conjuntos mecânicos e desenhos técnicos utilizando \textbf{SolidWorks}, \textbf{Autodesk Inventor} e \textbf{Fusion 360}.
    \item \textbf{Estudos de Viabilidade:} Análise técnica de projetos, levantamento de requisitos, atendimento ao cliente e elaboração de propostas comerciais.
\end{itemize}

% ==========================================
% HABILIDADES TÉCNICAS
% ==========================================
\section{Habilidades Técnicas}
\begin{itemize}[nosep, leftmargin=*]
    \item \textbf{CAD / Modelagem 3D:} SolidWorks, Autodesk Inventor, Fusion 360, Desenho Técnico Mecânico.
    \item \textbf{Automação \& CLPs:} TIA Portal (Linguagem Ladder), RobotStudio (ABB), SCADA (E3 Viewer), Sistemas Eletropneumáticos.
    \item \textbf{Programação / Firmware:} C/C++, Arduino (Microcontroladores), Lógica de Programação.
    \item \textbf{Simulação / FEA:} ANSYS (Estrutural), Altair Inspire (Otimização Topológica).
    \item \textbf{Fabricação & Impressão 3D:} Manufatura Aditiva, Aplicações e propriedades de materiais termoplásticos (PETG, TPU).
    \item \textbf{Idiomas / Outros:} Inglês (Leitura Técnica Autônoma de documentações/manuais), Microsoft Excel, Gestão de Projetos.
\end{itemize}

% ==========================================
% MÓDULOS DO CURSO TÉCNICO (SENAI - 908h)
% ==========================================
\section{Formação Técnica Detalhada (SENAI)}
\begin{itemize}[nosep, leftmargin=*]
    \item \textbf{Módulo I -- Eletroeletrônica e Controle (424h):} Acionamentos Eletroeletrônicos, Instrumentação e Controle de Processos Industriais, Sistemas Eletrohidráulicos e Eletropneumáticos, Sistemas Eletrônicos e Microcontrolados, Sistemas Lógicos Programáveis.
    \item \textbf{Módulo II -- Sistemas Automatizados (292h):} Gestão dos Processos de Implementação de Sistemas Automatizados, Comissionamento de Sistemas, Sistemas de Supervisão e Controle, Integração de Dispositivos Automatizados, Manutenção de Sistemas.
    \item \textbf{Módulo III -- Projetos (192h):} Projetos de Acionamentos Eletroeletrônicos, Projetos de Controle e Sistemas Automatizados, Projetos de Intertravamento de Segurança de Processos Industriais.
\end{itemize}

\end{document}
